
\documentclass[border=4pt]{standalone}
\usepackage{polyglossia}
\setdefaultlanguage{russian}
\setotherlanguage{english}
\usepackage{fontspec}
\IfFontExistsTF{CMU Serif}{\setmainfont{CMU Serif}}{\setmainfont{Liberation Serif}}
\IfFontExistsTF{CMU Sans Serif}{\setsansfont{CMU Sans Serif}}{\setsansfont{Liberation Sans}}
\IfFontExistsTF{CMU Typewriter Text}{\setmonofont{CMU Typewriter Text}}{\setmonofont{DejaVu Sans Mono}}
\usepackage{amsmath,amssymb,mathtools,bm}
\usepackage{xcolor}
\definecolor{accent}{HTML}{1F4E79}
\definecolor{accent2}{HTML}{B85450}
\definecolor{shade}{HTML}{F4F1EA}
\colorlet{prosvBlue}{accent}
\colorlet{prosvRed}{accent2}
\colorlet{prosvLight}{accent!12}
\colorlet{prosvCream}{shade}
\definecolor{prosvGold}{HTML}{E8B647}
\definecolor{prosvGreen}{HTML}{6A9F58}
\colorlet{prosvGray}{black!55}
\usepackage{tikz}
\usetikzlibrary{calc, arrows.meta, decorations.pathreplacing,
                positioning, shapes.geometric, shapes.misc, fit,
                backgrounds}
\usepackage{pgfplots}
\pgfplotsset{compat=1.18}
\newcommand{\R}{\mathbb{R}}
\newcommand{\N}{\mathbb{N}}
\newcommand{\Z}{\mathbb{Z}}
\newcommand{\eps}{\varepsilon}
\newcommand{\norm}[1]{\left\|#1\right\|}
\newcommand{\abs}[1]{\left|#1\right|}
\newcommand{\NOD}{\gcd}
\newcommand{\Eul}{\varphi}
\newcommand{\Zn}[1]{\mathbb{Z}_{#1}}
\newcommand{\modn}[1]{\,(\mathrm{mod}\,#1)}
\newcommand{\term}[1]{\textcolor{accent2}{\textbf{#1}}}
\newcommand{\engterm}[1]{\textit{#1}}
\DeclareMathOperator*{\argmin}{arg\,min}
\DeclareMathOperator*{\argmax}{arg\,max}
\begin{document}

\begin{tikzpicture}[scale=0.85, every node/.style={font=\small\ttfamily}]
  
  \foreach \i/\ch in {0/V,1/E,2/N,3/I,4/V,5/I,6/D,7/I,8/V,9/I,10/C,11/I}{
    \draw[prosvBlue] (\i*0.6, 1.2) rectangle (\i*0.6+0.55, 1.7);
    \node at (\i*0.6+0.275, 1.45) {\ch};
  }
  \node[left, font=\small\rmfamily] at (0,1.45) {Открытый текст:};
  
  
  \foreach \i in {0,...,11}{
    \draw[->,>=stealth, prosvGray] (\i*0.6+0.275, 1.15) -- (\i*0.6+0.275, 0.75);
  }
  
  
  \foreach \i/\ch in {0/Y,1/H,2/Q,3/L,4/Y,5/L,6/G,7/L,8/Y,9/L,10/F,11/L}{
    \draw[prosvRed] (\i*0.6, 0.2) rectangle (\i*0.6+0.55, 0.7);
    \node at (\i*0.6+0.275, 0.45) {\ch};
  }
  \node[left, font=\small\rmfamily] at (0,0.45) {Шифр-текст:};
  
  \node[right, font=\small\rmfamily] at (7.5,1.0) {\itshape\color{prosvGray}«пришёл, увидел, победил»};
\end{tikzpicture}
\end{document}
